\section*{第 \ref{ch:g11:seismic-structure} 章——地震波揭示地球内部}

\begin{solution}{exo:g11:seismic-structure:1}
\omterm{def:g11:seismic-structure:body-waves}{P 波}压缩并拉伸岩石，岩石沿传播方向运动；
\omterm{def:g11:seismic-structure:body-waves}{S 波}使岩石剪切，岩石的运动方向与传播方向垂直。\omterm{def:g11:seismic-structure:body-waves}{P 波}可以穿过固体和液体，
\omterm{def:g11:seismic-structure:body-waves}{S 波}只能穿过固体。\omterm{def:g11:seismic-structure:body-waves}{P 波}先被记录到，
因为它是两者中较快的一种。
\end{solution}

\begin{solution}{exo:g11:seismic-structure:2}
$40 \times 111.2 = \qty{4448}{km}$，约为 \qty{4450}{km}。$8900/111.2
\approx \ang{80}$。
\end{solution}

\begin{solution}{exo:g11:seismic-structure:3}
\qty{456}{s} 即 7~分 36~秒；\qty{822}{s} 即 13~分
42~秒。S--P 延迟为 $822 - 456 = \qty{366}{s}$，即 6~分
6~秒。
\end{solution}

\begin{solution}{exo:g11:seismic-structure:4}
$\sin i_c = 6.0/8.0 = 0.75$，因此 $i_c \approx \ang{48.6}$。
\end{solution}

\begin{solution}{exo:g11:seismic-structure:5}
外核厚度为 $5150 - 2891 = \qty{2259}{km}$。地核半径为
$6371 - 2891 = \qty{3480}{km}$，且 $3480/6371 \approx 0.55$：
略大于地球半径的一半。
\end{solution}

\begin{solution}{exo:g11:seismic-structure:6}
6~分 6~秒即 \qty{366}{s}：震中距为 \ang{40}，即
$40 \times 111.2 \approx \qty{4450}{km}$。单一延迟只能给出距离：
震中可以位于以台站为圆心、该距离为半径的圆周上的任一点。
要确定与三个台站距离均符合的唯一一点，
还需另外两个台站。
\end{solution}

\begin{solution}{exo:g11:seismic-structure:7}
$i_c \approx \ang{48.6}$，$\cos i_c \approx 0.661$。直达波：$200/6.0
\approx \qty{33.3}{s}$。\omterm{ex:g11:seismic-structure:sp}{首波}：$200/8.0 + 2 \times 35 \times 0.661/6.0
\approx 25.0 + 7.7 = \qty{32.7}{s}$。\omterm{ex:g11:seismic-structure:sp}{首波}先到达，约早
\qty{0.6}{s}。
\end{solution}

\begin{solution}{exo:g11:seismic-structure:8}
$\tan i_c \approx 1.134$：临界距离为 $2 \times 35 \times 1.134
\approx \qty{79}{km}$。\omterm{prop:g11:seismic-structure:crossover}{交叉距离}为 $2 \times 35 \times
\sqrt{14.0/2.0} = 70\sqrt 7 \approx \qty{185}{km}$。
\cref{exo:g11:seismic-structure:7} 中的台站位于 \qty{200}{km} 处，超过
\omterm{prop:g11:seismic-structure:crossover}{交叉距离}：\omterm{ex:g11:seismic-structure:sp}{首波}先到达，与前面的结果一致。
\end{solution}

\begin{solution}{exo:g11:seismic-structure:9}
$\sin i_c = 6.5/8.0 = 0.8125$，$i_c \approx \ang{54.3}$，且 $\cos i_c \approx
0.583$。\omterm{prop:g11:seismic-structure:head-wave}{截距时间}：$2 \times 6 \times 0.583/6.5 \approx \qty{1.1}{s}$。
\omterm{prop:g11:seismic-structure:crossover}{交叉距离}：$12\sqrt{14.5/1.5} \approx \qty{37}{km}$。在大陆下方，
\omterm{prop:g11:seismic-structure:crossover}{交叉距离}约为 \qty{165}{km}，因此相隔数十千米的台站能同时记录到两条线；
在海洋下方，超过约 \qty{37}{km} 后\omterm{ex:g11:seismic-structure:sp}{首波}已先到达，
而直达波线只有数十千米长，
台站之间须相隔数千米，
才能画出这条线。
\end{solution}

\begin{solution}{exo:g11:seismic-structure:10}
$\sin r = (8.06/13.7)\sin\ang{30} \approx 0.294$，因此 $r \approx \ang{17.1}$：
射线向法线方向弯折，即朝地心方向弯折。
因此射线沿较长的路径穿过地核，路径靠近其中部，
并从地球另一侧较远处射出：进入地核的 \omterm{def:g11:seismic-structure:body-waves}{P 波}射线，
其出射点距震中远远超过恰好擦过地核的射线。
\end{solution}

\begin{solution}{exo:g11:seismic-structure:11}
$12\,742/13.7 \approx \qty{930}{s}$，约 15.5~分，而波实际\omterm{def:g11:seismic-structure:travel-time}{走时}为
\qty{1210}{s}，约 20~分，晚了约 280~秒。即使按地幔中的最高速度计算，
波也应早得多到达，因此地球中心必定使其减速。
Oldham 由此推断，地球具有中心地核，
其中 \omterm{def:g11:seismic-structure:body-waves}{P 波}的传播速度比在地幔中慢。
\end{solution}

\begin{solution}{exo:g11:seismic-structure:12}
(1) 由命题可得 $t_i = 2h\cos i_c/v_1$，且 $\cos i_c = \sqrt{1 -
v_1^2/v_2^2} = \sqrt{v_2^2 - v_1^2}/v_2$。因此
\[
  t_i = \frac{2h\sqrt{v_2^2 - v_1^2}}{v_1 v_2},
  \qquad\text{故}\qquad
  h = \frac{t_i\,v_1 v_2}{2\sqrt{v_2^2 - v_1^2}}.
\]
(2) $h = 6.23 \times 6.2 \times 8.1/\bigl(2\sqrt{65.61 -
38.44}\bigr) \approx 312.9/10.43 \approx \qty{30.0}{km}$。
\end{solution}

\begin{solution}{exo:g11:seismic-structure:13}
(1) $(3480/6371)^3 \approx 0.163$：地核约占地球体积的 \qty{16}{\percent}。
(2) 平均密度是按体积加权的平均值：$5.52 = 0.163
\rho_c + 0.837 \times 4.45$，因此 $\rho_c \approx (5.52 - 3.72)/0.163 \approx
\qty{11.0}{g/cm^3}$。(3) 这一数值大于 \qty{9.9}{g/cm^3}，符合预期，
因为地核密度向中心增大（至 \qty{13.1}{g/cm^3}）；
它约为地壳密度
\qty{2.6}{g/cm^3} 的四倍。地核由密度远大于
地表岩石的物质构成。
\end{solution}

\begin{solution}{exo:g11:seismic-structure:14}
(1) $105 \times 111.2 \approx \qty{11700}{km}$，$143 \times 111.2 \approx
\qty{15900}{km}$。(2) $(\cos\ang{105} - \cos\ang{143})/2 \approx (-0.259 +
0.799)/2 \approx 0.27$：约占地表面积的 \qty{27}{\percent}，是围绕震中对面一点的一圈环带。
(3) \omterm{def:g11:seismic-structure:shadow-zone}{影区}是距震中的一段距离范围，
它随地震而移动。
每次地震都有各自的环带。
\end{solution}

\begin{solution}{exo:g11:seismic-structure:15}
(1) $h = 9.5 \times 6.3 \times 8.1/\bigl(2\sqrt{65.61 - 39.69}\bigr) \approx
484.8/10.18 \approx \qty{48}{km}$. (2) $\sin i_c = 6.3/8.1 \approx 0.778$,
$\cos i_c \approx 0.629$，且 $t_i = 2 \times 30 \times 0.629/6.3 \approx
\qty{6.0}{s}$。(3) 山脉之下的地壳厚约 \qty{48}{km}，
比大陆地壳平均厚度约 \qty{41}{km} 更厚，且接近
造山带的平均值：高山矗立于厚地壳之上。
\end{solution}

\begin{solution}{pb:g11:seismic-structure:1}
\textbf{1.} Pg 点位于一条直线上；Pn 点位于
另一条较平缓的直线上。

\textbf{2.} 地震发生在震中正下方很浅的深度：
距离为零的台站将立即记录到该波。

\textbf{3.} $(32.3 - 3.3)/(200 - 20) = 29.0/180 \approx \num{0.161}$~\unit{s}
每千米。

\textbf{4.} $v_1 = 180/29.0 \approx \qty{6.2}{km/s}$.

\textbf{5.} $3.3 + 280 \times 0.161 \approx \qty{48.4}{s}$.

\textbf{6.} $(43.2 - 18.6)/(300 - 100) = 24.6/200 = \num{0.123}$~\unit{s}
每千米，故 $v_2 = 200/24.6 \approx \qty{8.13}{km/s}$。

\textbf{7.} 接近 \qty{8.0}{km/s}：地壳之下是
最上部地幔的岩石。

\textbf{8.} $\sin i_c = 6.21/8.13 \approx 0.763$，故 $i_c \approx
\ang{49.8}$。

\textbf{9.} \omterm{ex:g11:seismic-structure:sp}{首波}以角度 $i_c$ 向下传播，并以相同角度
返回地表的首个能返回地表的点位于
到达\omterm{def:g11:seismic-structure:moho}{莫霍面}之点以外 $h\tan i_c$ 处，而该点距震中 $h\tan i_c$，
故距震中小于 $2h\tan i_c$ 的台站都无法接收它。

\textbf{10.} 地震仪记录到达它的每一个波，不只是
第一个波；Pn 波晚于 Pg 波到达，但在记录中
仍作为一个独立的波显示。

\textbf{11.} 在 \qty{160}{km} 处，Pg（\qty{25.8}{s}）先于 Pn
（\qty{26.0}{s}）到达；在 \qty{180}{km} 处，Pn（\qty{28.5}{s}）先于 Pg
（\qty{29.0}{s}）到达。这一变化发生在这两个台站之间。

\textbf{12.} Pg：$t = 3.3 + 0.1611\,(x - 20) \approx 0.08 + 0.1611\,x$。Pn：
$t = 18.6 + 0.123\,(x - 100) = 6.30 + 0.123\,x$（$x$ 以千米计，$t$ 以
秒计）。

\textbf{13.} 由 $0.08 + 0.1611\,x = 6.30 + 0.123\,x$ 得 $x \approx
6.22/0.0381 \approx \qty{163}{km}$。

\textbf{14.} Pn，时刻为 $6.30 + 0.123 \times 250 \approx \qty{37.1}{s}$。

\textbf{15.} $t_i \approx \qty{6.3}{s}$.

\textbf{16.} $\cos i_c \approx 0.646$，故 $h = 6.21 \times 6.3/(2 \times
0.646) \approx \qty{30.3}{km}$.

\textbf{17.} $x_c = 2 \times 30.3 \times \sqrt{14.34/1.92} \approx 60.6 \times
2.73 \approx \qty{165}{km}$，与第~13~题中的 \qty{163}{km}
相差不超过 \qty{2}{km}。

\textbf{18.} $\tan i_c \approx 1.18$: $2 \times 30.3 \times 1.18 \approx
\qty{72}{km}$。最近能记录到 Pn 的台站位于 \qty{80}{km} 处，而
\qty{60}{km} 处台站则不能，两者相符。

\textbf{19.} \omterm{def:g11:seismic-structure:travel-time}{走时}读数精确到 \qty{0.1}{s}，$t_i$ 上仅 \qty{0.1}{s} 的误差
便使 $h$ 改变约 \qty{0.5}{km}；速度
由斜率求得，同样不够确定，且地震并非恰好发生在
地表。地壳也并非完全均匀。数据所能给出的精度
最多为一到两千米。

\textbf{20.} \textbf{该地区\omterm{def:g11:seismic-structure:moho}{莫霍面}深约 \qty{30}{km}}：
此处地壳对大陆而言偏薄，大陆地壳平均厚约
\qty{41}{km}；\qty{30}{km} 为地壳经受
拉伸地区的典型厚度。
\end{solution}
